% ============================================================
%  BEAMER — T1L2: MASA MOLAR, FÓRMULA EMPÍRICA Y MOLECULAR
%  Curso: 4to Año — Química
%  Autor: Ing. Gabriel Astudillo
%  Fuente de la clase: t1l2 química.md (21 al 28 de octubre)
%
%  ESTRUCTURA
%   Portada + §Repaso exprés (3 láminas) +
%   §Ejercicios explicados (6 láminas, solución paso a paso) +
%   §Ejercicios en blanco (8 láminas, enunciado arriba y pizarra libre).
%
%  Los 8 ejercicios en blanco son DISTINTOS de los 60 de la guía:
%  sirven para resolver en vivo sin que el alumno tenga la respuesta.
%
%  Compilar desde esta carpeta:
%    latexmk -xelatex beamer_t1l2_formulas_quimicas.tex
% ============================================================

\documentclass[aspectratio=169]{beamer}
\usepackage{mathwizards-beamer}

\mwbeamer{T1L2}{Masa molar, fórmula empírica y molecular}{4to Año — Química}

% ── Comandos propios de esta presentación ────────────────────
% babel-spanish redefine \% y falla en modo matemático
% ("Incompatible glue units"): se usa \pc para el signo % en fórmulas.
\newcommand{\pc}{\text{\%}}

\begin{document}

\begin{frame}[plain]
  \titlepage
\end{frame}

% ════════════════════════════════════════════════════════════
\section{Repaso exprés}
% ════════════════════════════════════════════════════════════

\begin{frame}{Masa molar: sumar con subíndices}
  \begin{block}{La fórmula}
    \[
      M(\text{compuesto}) = \sum \bigl(\text{n.º de átomos}\bigr)\times \bigl(M.A.\bigr)
      \qquad [\ \text{g/mol}\ ]
    \]
  \end{block}
  \vspace{2pt}
  \begin{mwpasos}[Cómo se lee un subíndice]
    \begin{itemize}
      \item $\text{H}_2\text{O}$: dos H y un O $\Rightarrow 2(1)+16 = 18$.
      \item El subíndice afecta a \textbf{todo} el paréntesis:
            $\text{Ca(OH)}_2 \Rightarrow 40 + 2(16) + 2(1) = 74$.
      \item $\text{Al}_2(\text{SO}_4)_3 \Rightarrow$ Al: 2, S: 3, O: $3\times4 = 12$.
    \end{itemize}
  \end{mwpasos}
  \vspace{2pt}
  \begin{mwextra}[Masas atómicas que usaremos]
    \centering
    \small
    H $=1$ \quad C $=12$ \quad N $=14$ \quad O $=16$ \quad Na $=23$ \quad Mg $=24{,}3$\\[2pt]
    Al $=27$ \quad P $=31$ \quad S $=32$ \quad Cl $=35{,}5$ \quad K $=39$ \quad Ca $=40$ \quad Fe $=56$
  \end{mwextra}
\end{frame}

\begin{frame}{Composición porcentual y la ruta completa}
  \begin{block}{Dos fórmulas, un camino}
    \[
      \pc\,\text{elemento} = \frac{\bigl(\text{n.º de átomos}\bigr)\times \bigl(M.A.\bigr)}{M(\text{compuesto})}\times 100
      \qquad
      n = \frac{M(\text{compuesto})}{M(\text{F.E.})}
    \]
  \end{block}
  \vspace{2pt}
  \begin{mwpasos}[Ruta del porcentaje a la fórmula molecular]
    \begin{enumerate}
      \item Suponer $100$ g de muestra: cada $\pc$ pasa a gramos.
      \item Gramos a moles: $n = \text{g}/M.A.$
      \item Dividir entre el \textbf{menor} número de moles.
      \item Convertir a enteros con la tabla de factores.
      \item Calcular $n$ y multiplicar la F.E.: $\text{F.M.} = n\times\text{F.E.}$
    \end{enumerate}
  \end{mwpasos}
\end{frame}

\begin{frame}{Multiplicadores y errores clásicos}
  \begin{block}{Cuando la relación no da entero}
    \centering
    \small
    Termina en $\sim0{,}50 \Rightarrow \times2$
    \qquad
    $\sim0{,}33$ o $\sim0{,}66 \Rightarrow \times3$
    \qquad
    $\sim0{,}25$ o $\sim0{,}75 \Rightarrow \times4$
  \end{block}
  \vspace{2pt}
  \begin{mwextra}[Errores que se pagan caro]
    \begin{itemize}
      \item Dividir entre el \emph{primer} valor en vez del \emph{menor}.
      \item Redondear $1{,}33$ a $1$: es $4/3$, ¡obliga a multiplicar por $3$!
      \item Multiplicar por el factor a un solo elemento.
      \item Confundir $\text{CH}_2\text{O}$ (F.E.) con $\text{C}_6\text{H}_{12}\text{O}_6$ (F.M.).
    \end{itemize}
  \end{mwextra}
\end{frame}

% ════════════════════════════════════════════════════════════
\section{Ejercicios explicados}
% ════════════════════════════════════════════════════════════

\begin{frame}{Ejercicio 1 — Masa molar con paréntesis}
  \mwtopanchoroff
  \begin{mwejercicio}[Ejercicio 1 — Sal con paréntesis \quad \medio]
    Calcule la masa molar del fosfato de calcio, $\text{Ca}_3(\text{PO}_4)_2$.
  \end{mwejercicio}
  \vspace{2pt}
  \begin{mwpasos}[Resolución]
    \begin{itemize}
      \item Contar átomos: el subíndice $2$ multiplica todo el paréntesis.
            \[
              \text{Ca}:3 \qquad \text{P}:2\times1 = 2 \qquad \text{O}:2\times4 = 8
            \]
      \item Multiplicar y sumar:
            \[
              M = 3(40) + 2(31) + 8(16) = 120 + 62 + 128 = \boxed{310\ \text{g/mol}}
            \]
    \end{itemize}
  \end{mwpasos}
  \mwtopanchoron
\end{frame}

\begin{frame}{Ejercicio 2 — Masa molar de una sal compleja}
  \mwtopanchoroff
  \begin{mwejercicio}[Ejercicio 2 — Sulfato de hierro(III) \quad \dificil]
    Calcule la masa molar del sulfato de hierro(III), $\text{Fe}_2(\text{SO}_4)_3$.
  \end{mwejercicio}
  \vspace{2pt}
  \begin{mwpasos}[Resolución]
    \begin{itemize}
      \item Contar átomos: Fe $=2$, S $=3$, O $=3\times4 = 12$.
      \item Multiplicar y sumar:
            \[
              M = 2(56) + 3(32) + 12(16) = 112 + 96 + 192 = \boxed{400\ \text{g/mol}}
            \]
      \item \textbf{Chequeo rápido:} el O siempre es el que más aporta
            ($192$ de $400$, casi la mitad).
    \end{itemize}
  \end{mwpasos}
  \mwtopanchoron
\end{frame}

\begin{frame}{Ejercicio 3 — Composición porcentual}
  \mwtopanchoroff
  \begin{mwejercicio}[Ejercicio 3 — Óxido de aluminio \quad \dificil]
    Determine la composición porcentual del óxido de aluminio, $\text{Al}_2\text{O}_3$.
  \end{mwejercicio}
  \vspace{2pt}
  \begin{mwpasos}[Resolución]
    \begin{itemize}
      \item Masa molar: $M = 2(27) + 3(16) = 54 + 48 = 102$ g/mol.
      \item Porcentajes:
            \[
              \pc\,\text{Al} = \frac{54}{102}\times100 = 52{,}94\pc
              \qquad
              \pc\,\text{O} = \frac{48}{102}\times100 = 47{,}06\pc
            \]
      \item Verificación: $52{,}94 + 47{,}06 = 100\pc$ \checkmark
    \end{itemize}
  \end{mwpasos}
  \mwtopanchoron
\end{frame}

\begin{frame}{Ejercicio 4 — Fórmula empírica desde porcentajes}
  \mwtopanchoroff
  \begin{mwejercicio}[Ejercicio 4 — Clorato de potasio \quad \dificil]
    Un compuesto contiene $31{,}84\pc$ de K, $28{,}98\pc$ de Cl y $39{,}18\pc$ de O.
    Determine su fórmula empírica.
  \end{mwejercicio}
  \vspace{2pt}
  \begin{mwpasos}[Resolución — base de cálculo $100$ g]
    \centering
    \small
    \setlength{\tabcolsep}{4pt}
    \begin{tabular}{lcccc}
      \toprule
      Elemento & Gramos & Moles & Relación & Subíndice \\
      \midrule
      K  & $31{,}84$ & $31{,}84/39 = 0{,}816$    & $0{,}816/0{,}816 = 1$ & $1$ \\
      Cl & $28{,}98$ & $28{,}98/35{,}5 = 0{,}816$ & $0{,}816/0{,}816 = 1$ & $1$ \\
      O  & $39{,}18$ & $39{,}18/16 = 2{,}449$     & $2{,}449/0{,}816 = 3$ & $3$ \\
      \bottomrule
    \end{tabular}
    \par\vspace{4pt}
    $\text{F.E.} = \boxed{\text{KClO}_3}$ \quad (clorato de potasio)
  \end{mwpasos}
  \mwtopanchoron
\end{frame}

\begin{frame}{Ejercicio 5 — De la fórmula empírica a la molecular}
  \mwtopanchoroff
  \begin{mwejercicio}[Ejercicio 5 — Hidrocarburo \quad \medio]
    La fórmula empírica de un hidrocarburo es $\text{CH}_2$ y su masa molar es
    $70$ g/mol. Determine su fórmula molecular.
  \end{mwejercicio}
  \vspace{2pt}
  \begin{mwpasos}[Resolución]
    \begin{itemize}
      \item Masa molar de la fórmula empírica: $M(\text{CH}_2) = 12 + 2(1) = 14$ g/mol.
      \item Factor: $n = 70/14 = 5$.
      \item Multiplicar: $\text{F.M.} = 5\times\text{CH}_2 = \boxed{\text{C}_5\text{H}_{10}}$.
      \item Se trata de un alqueno (penteno).
    \end{itemize}
  \end{mwpasos}
  \mwtopanchoron
\end{frame}

\begin{frame}{Ejercicio 6 — Ruta completa}
  \mwtopanchoroff
  \begin{mwejercicio}[Ejercicio 6 — F.E. y F.M. desde el porcentaje \quad \muyDificil]
    Un compuesto contiene $40{,}00\pc$ de C, $6{,}67\pc$ de H y $53{,}33\pc$ de O.
    Su masa molar es $180$ g/mol. Determine la F.E. y la F.M.
  \end{mwejercicio}
  \vspace{2pt}
  \begin{mwpasos}[Resolución]
    \small
    \begin{itemize}
      \item Moles (base $100$ g): C $=40/12 = 3{,}33$; H $=6{,}67/1 = 6{,}67$; O $=53{,}33/16 = 3{,}33$.
      \item Dividir entre $3{,}33$: C : H : O $= 1 : 2 : 1 \Rightarrow \text{F.E.} = \text{CH}_2\text{O}$.
      \item $M(\text{CH}_2\text{O}) = 12 + 2(1) + 16 = 30$ g/mol $\Rightarrow n = 180/30 = 6$.
      \item $\text{F.M.} = 6\times\text{CH}_2\text{O} = \boxed{\text{C}_6\text{H}_{12}\text{O}_6}$ \ (glucosa).
    \end{itemize}
  \end{mwpasos}
  \mwtopanchoron
\end{frame}

% ════════════════════════════════════════════════════════════
\section{Ejercicios en blanco}
% ════════════════════════════════════════════════════════════

\begin{frame}{Ejercicio 7}
  \begin{mwejercicio}[Ejercicio 7 — Masa molar \quad \medio]
    Calcule la masa molar del nitrato de amonio, $\text{NH}_4\text{NO}_3$,
    un fertilizante muy usado.
  \end{mwejercicio}
\end{frame}

\begin{frame}{Ejercicio 8}
  \begin{mwejercicio}[Ejercicio 8 — Sal hidratada \quad \dificil]
    Calcule la masa molar del sulfato de cobre pentahidratado,
    $\text{CuSO}_4\cdot5\,\text{H}_2\text{O}$.
    \par\smallskip
    \small
    \textbf{Sugerencia:} el $5\,\text{H}_2\text{O}$ se suma como $5\times18$.
  \end{mwejercicio}
\end{frame}

\begin{frame}{Ejercicio 9}
  \begin{mwejercicio}[Ejercicio 9 — Composición porcentual \quad \dificil]
    Determine la composición porcentual completa del dicromato de potasio,
    $\text{K}_2\text{Cr}_2\text{O}_7$, y verifique que los porcentajes sumen $100\pc$.
  \end{mwejercicio}
\end{frame}

\begin{frame}{Ejercicio 10}
  \begin{mwejercicio}[Ejercicio 10 — Fórmula empírica directa \quad \medio]
    Un compuesto contiene $46{,}67\pc$ de nitrógeno y $53{,}33\pc$ de oxígeno.
    Determine su fórmula empírica.
  \end{mwejercicio}
\end{frame}

\begin{frame}{Ejercicio 11}
  \begin{mwejercicio}[Ejercicio 11 — Fórmula empírica con factor \quad \muyDificil]
    Un óxido de nitrógeno contiene $25{,}93\pc$ de nitrógeno y $74{,}07\pc$ de
    oxígeno. Determine su fórmula empírica.
    \par\smallskip
    \small
    \textbf{Sugerencia:} la relación molar no es entera; revise la tabla de factores.
  \end{mwejercicio}
\end{frame}

\begin{frame}{Ejercicio 12}
  \begin{mwejercicio}[Ejercicio 12 — Fórmula molecular \quad \dificil]
    La fórmula empírica de un hidrocarburo es $\text{CH}_2$ y su masa molar es
    $56$ g/mol. Determine su fórmula molecular e indique cuántos carbonos tiene.
  \end{mwejercicio}
\end{frame}

\begin{frame}{Ejercicio 13}
  \begin{mwejercicio}[Ejercicio 13 — Ruta completa \quad \muyDificil]
    Un compuesto contiene $92{,}26\pc$ de carbono y $7{,}74\pc$ de hidrógeno.
    Su masa molar es $78$ g/mol. Determine la fórmula empírica y la molecular.
  \end{mwejercicio}
\end{frame}

\begin{frame}{Ejercicio 14}
  \begin{mwejercicio}[Ejercicio 14 — Masa de un elemento \quad \muyDificil]
    ¿Cuántos gramos de nitrógeno hay en $250$ g de nitrato de amonio,
    $\text{NH}_4\text{NO}_3$?
    \par\smallskip
    \small
    \textbf{Sugerencia:} primero el $\pc$ de N en el compuesto; después, la masa en la muestra.
  \end{mwejercicio}
\end{frame}

\end{document}
